\section{Introduction}
\label{sec:introduction}

Mathematical reasoning requires language models to produce not merely
correct final answers but logically sound, step-by-step derivations in
which each conclusion follows from its
premises~\citep{ElKishky2024OpenAIOS,Chen2025TowardsRE}. Reinforcement
learning with verifiable rewards
(RLVR)~\citep{deepseekr1,shao2024deepseekmath} has become a standard
approach for improving such reasoning: a programmatic check verifies
whether the model's answer matches the ground truth, and the resulting
binary signal drives policy optimization through Group Relative Policy
Optimization (GRPO). Because the reward depends only on the final answer,
however, RLVR provides no gradient for improving the quality of the
reasoning trace itself.

This outcome-only supervision creates three systematic failure modes.
First, models can reach correct answers without constructing logically
sound derivations, a phenomenon sometimes called outcome
hacking~\citep{Denison2024SycophancyTS}. Second, reasoning traces often
contain orphan conclusions that lack supporting premises, circular
references between steps, and inconsistent logical direction, because
inter-step dependencies are never supervised. Third, without any length
or coherence signal, models tend to produce verbose, repetitive chains
that waste inference budget and occasionally introduce
self-contradictions~\citep{reasoning_compression_2025}.

Process reward models (PRMs)~\citep{lightman2023lets,wang2024math}
address some of these issues by assigning per-step scores, but they
require either expensive human annotations or learned verifiers whose
accuracy is itself uncertain. We observe that a complementary source of
process supervision is already present in the reasoning trace: the
\emph{topological structure} of inter-step dependencies. Mathematical
derivations form a directed acyclic graph (DAG) in which each step
depends on a subset of previous steps. Structural properties of this
graph, such as acyclicity, connectivity, and directional consistency, are
checkable by simple algorithms and provide verifiable process rewards
without any learned component.

We propose \textbf{TopoPRM}, a composite reward framework that
supplements the outcome reward with two topology-aware process signals.
Given a model-generated reasoning trace, a rule-based parser segments it
into steps and detects inter-step dependencies through shared
mathematical expressions, producing the implicit reasoning DAG. A
\emph{topological structure reward} evaluates whether this DAG satisfies
structural soundness properties. A \emph{continuity reward} checks
whether each step can be derived from preceding context, penalizing
unsupported jumps. Both rewards are computed by deterministic programs
and stay within the RLVR paradigm.

Combining multiple reward components introduces a new risk: the model
may optimize auxiliary process rewards at the expense of answer
correctness. To prevent this, we introduce \emph{Stratified Clipping
Advantage Estimation} (SCAE), which partitions each group of sampled
completions into correct and incorrect sets before computing advantages.
SCAE ensures that correct completions always receive positive advantages
regardless of their structural scores, while incorrect completions
receive negative advantages regardless of how well-structured they are.
After GRPO training, a topology-guided self-distillation stage compresses
the teacher's DAG-verified reasoning behaviour into a compact ($\leq$4B)
student. The goal of this stage is explicitly \emph{not} to match teacher
accuracy, but to retain the teacher's structural reasoning profile
(acyclic\%, no-orphan\%, edge-keep\%) while halving the inference-time
generation budget, at the cost of a mild 1--3~pp absolute accuracy drop
on math benchmarks. In summary, our contributions are fourfold:
\begin{enumerate}[leftmargin=*,itemsep=2pt]
\item \textbf{Verifiable process rewards.} We formulate process supervision for RLVR as two deterministic, complementary signals over an implicit reasoning DAG: a topological structure reward for global graph validity and a continuity reward for local step traceability. Neither requires human step annotations nor learned reward models.

\item \textbf{Accuracy-first reward shaping.} We introduce Stratified Clipping Advantage Estimation (SCAE), which partitions sampled completions by correctness before computing advantages, preventing process rewards from favouring structurally valid but factually wrong traces.

\item \textbf{End-to-end training recipe with topology-guided compression.} We instantiate TopoPRM as plug-in reward modules compatible with standard GRPO stacks, and add a \emph{topology-verified self-distillation} stage (TVSD) that reuses the same deterministic diagnostics to densify the scalar outcome signal into per-token supervision. TVSD produces a compact $\leq$4B student that retains the teacher's DAG structural profile while halving the generation budget, without any external teacher or human step-level annotation.

\item \textbf{Comprehensive empirical validation.} On 7{,}595 Chinese mathematical critique problems, TopoPRM improves critique accuracy by $+13.4$ points over outcome-only GRPO while reducing response length by 11\%. On public math benchmarks, the compact TVSD 4B student halves the teacher's output tokens and retains $\geq$95\% of its structural-retention metrics, at the cost of 1--3~pp absolute accuracy drop---establishing a useful Pareto point on the accuracy--tokens frontier.
\end{enumerate}